% NAECOMPANION1 2026-10-18 16:55 Asia/Seoul
\documentclass[11pt]{article}
\usepackage[T1]{fontenc}
\usepackage{lmodern,amsmath,amssymb,amsthm,booktabs,hyperref}
\usepackage[margin=.9in]{geometry}
\newtheorem{theorem}{Theorem}
\newtheorem{lemma}[theorem]{Lemma}
\newtheorem{corollary}[theorem]{Corollary}
\newcommand{\EXACT}{\textbf{PROVED/EXACT}}
\newcommand{\KNOWN}{\textbf{KNOWN THEOREM}}
\newcommand{\COMP}{\textbf{COMPUTATIONAL}}
\newcommand{\CAND}{\textbf{CANDIDATE}}
\newcommand{\OPEN}{\textbf{OPEN}}
\newcommand{\RET}{\textbf{FALSIFIED/RETIRED}}
\title{NAECOMPANION1:\\Exact SAT/UNSAT Gauge Companions and Constant-Spread DNF Transfer}
\author{P versus NP project}
\date{18 October 2026}
\begin{document}\maketitle

\noindent\textbf{Status.}
The exact-density SAT/UNSAT companion construction, exact local-moment identity, spread-core peeling theorem, and bounded-width DNF transfer theorem are \EXACT at their declared distributional models, pending independent review. The finite Fano/cycle/decomposition audit is \COMP. Ellis--Linial high-girth hypergraphs and CKR robust-sunflower approximation are \KNOWN. No circuit-to-DNF theorem with the required uniform width coefficient is proved; no unrestricted monotone/set-circuit lower bound or $P\ne NP$ proof is obtained.

\section{IDEAS PROPOSED BEFORE LITERATURE}
\paragraph{A: same-scaffold SAT companion.}
Use the same high-girth support, random gauge, and exact-density transparent completion for two base signings, one UNSAT and one SAT. Prediction: every clause-coordinate inclusion moment below the support girth is exactly identical, not merely asymptotically close.

\paragraph{B: spread-core peeling.}
Repeatedly locate one dense core whose residual link is spread, remove the represented terms, and continue. Prediction: a width-$W$ family decomposes into at most $\rho^W(1+\log L)$ spread clusters, where $L$ is the term count.

\paragraph{C: constant-density gauge transfer.}
On a fixed high-girth scaffold, each support edge selects one of four NAE sign classes. Below half the girth, first and second term moments are exactly those of independent $1/4$-biased coordinates. Prediction: arbitrary width-$W$ DNFs transfer between the SAT and UNSAT gauge companions with only $\exp(O(W))$ loss.

\paragraph{D: ordinary quiet planting.}
Search planted-CSP literature for a ready-made SAT companion statistically identical to an UNSAT distribution. Falsifier: the available theorems are model- or algorithm-specific and do not provide exact all-pattern equality or an unrestricted monotone-circuit conclusion.

\section{LITERATURE CHECK AND REVISIONS}
Ellis--Linial construct infinite regular uniform hypergraph sequences with logarithmic Berge girth; the project uses their 3-uniform 9-regular sequence with
\[
g(H_v)\ge \gamma\log v-O(1),\qquad \gamma=\frac{2}{3\log 11}.
\]
The core degree nine is also sufficient for the elementary first-moment existence of an UNSAT signed NAE labeling.

CKR use robust sunflowers to approximate monotone circuits by bounded-width monotone functions and obtain strong monotone lower bounds. Their stock gate-by-gate approximation still incurs circuit-size-times-local-error accumulation, so it does not supply the final exponent-uniform width bound needed here.

Recent quiet-planting work gives planted SAT distributions indistinguishable from uniform for restricted algorithmic models, including statistical-query analyses. It does not give the exact companion identity below, nor a lower bound for unrestricted monotone circuits. Candidate D is therefore modified rather than adopted.

\section{CONSTRUCTION OR PROOF ATTEMPT}
\subsection{An exact-density transparent completion}
Let $H_v=(V,E)$ be a 3-uniform 9-regular hypergraph on $v$ vertices. Then $|E|=3v$. For every $u\in V$, introduce fresh vertices $a_u,b_u$ and one leaf edge
\[
L_u=\{u,a_u,b_u\}.
\]
Let $\widehat H_v$ be the augmented support. It has
\[
q=3v\quad\text{vertices},\qquad t=4v=\frac{4q}{3}\quad\text{edges}.
\]
A leaf edge lies in no Berge cycle, so
\[
g(\widehat H_v)=g(H_v).
\]

A signed NAE predicate on an edge $e$ is represented by
\[
Q_e=\mathbb F_2^3/\langle(1,1,1)\rangle,
\]
which has four sign classes. Choose:
\begin{itemize}
\item an UNSAT signing $\sigma^+$ on the core $H_v$, whose existence follows from
\[
\mathbb E[\#\mathrm{solutions}]
=2^v(3/4)^{3v}=(27/32)^v<1;
\]
\item a SAT signing $\sigma^-$ on the core, for example one sign class per edge satisfied by the all-zero assignment;
\item the same all-zero-satisfiable sign class on every leaf edge for both formulas.
\end{itemize}

For a uniform gauge vector $x\in\mathbb F_2^{V(\widehat H_v)}$, add $x|_e$ modulo $\langle111\rangle$ to the sign class on every edge. Finally apply a uniform permutation of the $q$ variable names. Let the resulting distributions be $\mu_v^+$ and $\mu_v^-$. Variable complementation and renaming preserve satisfiability, so
\[
\mu_v^+\text{ is supported on exact-density UNSAT formulas},
\]
while
\[
\mu_v^-\text{ is supported on exact-density SAT formulas}.
\]
The leaf clauses are transparent: once the core assignment is fixed, their two private variables can be chosen to satisfy them; for the displayed base signing, the all-zero assignment already does so.

\subsection{Exact local indistinguishability}
\begin{theorem}[Exact SAT/UNSAT companion moments]
For every set $A$ of signed NAE clause coordinates with
\[
|A|<g(\widehat H_v),
\]
one has
\[
\boxed{\Pr_{\mu_v^+}[A\subseteq F]=\Pr_{\mu_v^-}[A\subseteq F].}
\]
\end{theorem}
\begin{proof}
Condition on the variable permutation. If a requested coordinate is not supported by $\widehat H_v$, or if $A$ requests two different sign classes on one support edge, both probabilities are zero. Otherwise the supports requested by $A$ form an edge set $J$ of size less than the girth, hence a Berge hyperforest.

For a hyperforest, the gauge map
\[
\delta_J:\mathbb F_2^{V(J)}\to\bigoplus_{e\in J}Q_e
\]
is onto. Thus the transformed sign classes on $J$ are mutually independent and uniform, irrespective of the base signing. The requested pattern has probability $4^{-|J|}$ under both distributions. Averaging over the common permutation proves the claim.
\end{proof}

This gives an exact SAT-supported negative distribution. It removes the rare-UNSAT contamination present when the uniform exact layer is used directly.

\subsection{Spread-core peeling}
Let $\mathcal F$ be a $w$-uniform family of distinct terms and $L=|\mathcal F|$. A residual family $\mathcal G$ is $\rho$-spread if
\[
d_{\mathcal G}(A)\le |\mathcal G|\rho^{-|A|}
\]
for every $A$.

\begin{lemma}[One spread cluster]
For every nonempty $w$-uniform family $\mathcal H$, there is a core $C$ such that
\[
|\{T\in\mathcal H:C\subseteq T\}|\ge |\mathcal H|\rho^{-|C|}
\]
and the link
\[
\mathcal G_C=\{T\setminus C:T\in\mathcal H,\ C\subseteq T\}
\]
is $\rho$-spread.
\end{lemma}
\begin{proof}
Start with $C=\varnothing$. If the current link is not spread, choose a nonempty violating set $A$, add it to $C$, and restrict to terms containing it. Each step retains more than a $\rho^{-|A|}$ fraction. Core size strictly increases and is at most $w$, so the process terminates with the stated density and spreadness.
\end{proof}

\begin{theorem}[Peeling decomposition]
Every $w$-uniform family of $L$ terms can be partitioned into at most
\[
\boxed{J_w\le \rho^w(1+\log L)}
\]
spread clusters.
\end{theorem}
\begin{proof}
Apply the preceding lemma to the current remaining family, output the selected cluster, and remove all its represented terms. Every step removes at least a $\rho^{-w}$ fraction of the current terms. Thus after $s$ steps the remaining count is at most
\[
L(1-\rho^{-w})^s\le L\exp(-s\rho^{-w}).
\]
The displayed bound makes this quantity smaller than one. The removed term sets are disjoint, so the clusters form a partition.
\end{proof}

For a mixed-width family of width at most $W$, decompose each width separately. On a support with $N_s$ sign coordinates,
\[
L_w\le\binom{N_s}{w},
\]
so the total cluster count obeys
\[
J\le\sum_{w=0}^W\rho^w(1+\log L_w)
\le W\rho^W\left(1+W\log\frac{eN_s}{W}\right).
\]

\subsection{Companion transfer for arbitrary bounded-width DNFs}
Fix the support permutation and set unsupported coordinates to zero. Delete impossible terms requesting two sign classes on one edge. Let $D$ be the resulting monotone DNF of width at most $W$, and assume
\[
2W<g(\widehat H_v).
\]

Consider one spread cluster with core size $c$, residual width $r$, and residual family $\mathcal G$. Let
\[
R=\sum_{T\in\mathcal G}1[T\text{ occurs}],\qquad
\lambda=|\mathcal G|4^{-r}.
\]
The union of the core and any two residual terms uses fewer than $g$ support edges. It is therefore a forest. Under either companion distribution, conditioned on the core pattern,
\[
\mathbb E R=\lambda.
\]
Using spreadness exactly as in the standard second-moment calculation,
\[
\mathbb E(R)_2\le \lambda^2\exp(4r/\rho).
\]
Consequently, under the SAT companion,
\[
\eta_-:=\Pr[R>0\mid C]
\ge\frac{\lambda}{1+e^{4r/\rho}\lambda}.
\]
Under the UNSAT companion, a union bound gives
\[
\Pr[R>0\mid C]\le\lambda.
\]
If $\eta_-\le(2e^{4r/\rho})^{-1}$, rearrangement gives $\lambda\le2\eta_-$. Otherwise the trivial bound one is at most $2e^{4r/\rho}\eta_-$. Since the core itself has the same probability $4^{-c}$ under both companions, each cluster satisfies
\[
\Pr_{\mu_v^+}[\text{cluster}]
\le2e^{4W/\rho}
\Pr_{\mu_v^-}[\text{cluster}].
\]

\begin{theorem}[Term-count-free companion transfer]
For every monotone DNF $D$ of width at most $W$ with $2W<g(\widehat H_v)$,
\[
\boxed{
\Pr_{\mu_v^+}[D=1]
\le
2e^{4W/\rho}J\,
\Pr_{\mu_v^-}[D=1],
}
\]
where
\[
J\le W\rho^W\left(1+W\log\frac{eN_s}{W}\right),
\qquad N_s=4t=\Theta(q).
\]
The bound is independent of the number of original DNF terms except for the logarithm already absorbed in $J$.
\end{theorem}
\begin{proof}
Apply the cluster comparison and sum over the partition. Every cluster event is contained in $D$, so
\[
\sum_j\Pr_{\mu_v^-}[E_j]\le J\Pr_{\mu_v^-}[D].
\]
The bound is uniform in the support permutation; average over it.
\end{proof}

\subsection{The logarithmic-width cone}
For the Ellis--Linial sequence,
\[
g(\widehat H_v)\ge\gamma\log q-O(1),
\qquad \gamma=\frac{2}{3\log11}=0.27802159\ldots.
\]
Put
\[
W=\kappa\log q,
\qquad 0<\kappa<\gamma/2=0.13901079\ldots.
\]
Then the transfer loss has exponent
\[
\log\left(2e^{4W/\rho}J\right)
\le W\left(\log\rho+\frac4\rho\right)+o(\log q).
\]
The optimal fixed choice is $\rho=4$, yielding
\[
\log\rho+4/\rho=1+\log4=2.38629436\ldots.
\]
Therefore
\[
\boxed{
\Pr_{\mu_v^+}[D=1]
\le
q^{(1+\log4)\kappa+o(1)}
\Pr_{\mu_v^-}[D=1].
}
\]
At the edge of the available width window the exponent loss is below $0.33173$.

\begin{corollary}[Inverse-polynomial invisibility]
If for some fixed $\varepsilon>0$,
\[
\Pr_{\mu_v^-}[D=1]
\le q^{-((1+\log4)\kappa+\varepsilon)},
\]
then
\[
\Pr_{\mu_v^+}[D=1]=o(1).
\]
\end{corollary}

This replaces the previous quasi-polynomially small negative-error requirement by an inverse-polynomial requirement.

\subsection{Hostile control and remaining circuit theorem}
For a fixed labeled support, two distinct gauge cosets are separated by a dual cycle-parity check. Such a parity can have a linear-size ordinary circuit. Therefore exact local moment matching alone is not a general circuit lower bound. Random variable permutation removes the most trivial hardwired check but does not prove that the hidden support or its syndrome is computationally inaccessible.

A sufficient final theorem is now sharply stated. For every fixed circuit-size exponent $K$, produce a one-sided monotone DNF approximation of width
\[
W_K=\kappa_K\log q,
\qquad \kappa_K<\gamma/2,
\]
whose miss probability on $\mu_v^+$ is $o(1)$ and whose acceptance probability on $\mu_v^-$ is at most
\[
q^{-B_K},
\qquad B_K>(1+\log4)\kappa_K.
\]
The companion-transfer theorem would then give a contradiction. Standard gate-wise CKR accounting generally makes the width coefficient grow with $K$, so it does not establish this exponent-uniform window. The remaining problem is a global circuit-origin approximation or hidden-syndrome theorem, not construction of a negative distribution.

\section{TESTS AND OUTCOME}
The dependency-free audit used seed 20261018.

On the Fano-plane support, the all-positive NAE signing had zero satisfying assignments, while a one-negation-per-edge signing had 24. Their gauge orbits had 64 distinct formulas each and were disjoint. All 28 singleton and 378 pair inclusion moments agreed exactly. A three-edge cycle pattern distinguished the cosets, with probabilities $1/32$ and $0$, confirming that equality ends precisely at the cycle scale.

A loose five-edge NAE cycle was used to test the DNF transfer at width two. Two hundred random term families, containing up to 80 terms, were decomposed by spread-core peeling. Reconstruction, spreadness, the cluster-count bound, and both directions of the companion-transfer inequality were checked. The maximum observed acceptance ratio was $1.01841$, while the certified bounds were larger as expected.

Including density arithmetic, the audit performed 2,193 declared assertions with zero failures. Script SHA-256:
\[
\texttt{8f4a4782b08a06ea8da06280fe2f8965edc1c99a6781a9a80bb0a2d5db388fcd}.
\]
Raw-result SHA-256:
\[
\texttt{f83d82c27d4bc9599dee7bd7572256329c99f6db10f974907c300e2681930302}.
\]

\paragraph{Outcome.}
Candidate A is \EXACT: there is now a SAT-supported exact-density companion whose local clause moments match the UNSAT distribution exactly through logarithmic order. Candidate B is \EXACT. Candidate C is \EXACT and gives only polynomial loss at logarithmic width. Candidate D is not needed for the construction and remains a restricted-model literature comparison.

The negative-distribution obstacle is therefore closed at the bounded-width DNF level. The smallest remaining proof obligation is an exponent-uniform circuit-to-DNF or hidden-syndrome theorem. It must avoid a gate-by-gate factor proportional to the hypothetical polynomial circuit exponent.

\paragraph{Next target.}
Audit random-restriction and switching approaches directly on the companion pair, seeking a global approximation whose width coefficient is independent of circuit size exponent. A second route is to randomize the support scaffold itself and prove that recovering a distinguishing cycle syndrome requires large monotone circuits. Fixed-support parity and general-circuit canonicalization must remain hostile controls.

Root \texttt{main.tex} remains truncated and was not modified.

\begin{thebibliography}{6}
\bibitem{EL14} D. Ellis and N. Linial, \emph{On Regular Hypergraphs of High Girth}, Electron. J. Combin. 21(1):P1.54, 2014.
\bibitem{CKR22} B. P. Cavalar, M. Kumar, B. Rossman, \emph{Monotone Circuit Lower Bounds from Robust Sunflowers}, Algorithmica 84:3655--3685, 2022.
\bibitem{Quiet26} A. Ahmadi, K. Banihashem, I. Gholami, M. T. Hajiaghayi, J. Olkowski, \emph{Quiet Planting for $k$-SAT, Multiple Solutions of Arbitrary Geometry}, arXiv:2606.15979, 2026.
\bibitem{Kun13} G. Kun, \emph{Constraints, MMSNP and Expander Relational Structures}, Combinatorica 33:475--486, 2013.
\bibitem{ProjectGauge} P versus NP project, \emph{NAEGAUGEHIGHGIRTH1}, 17 October 2026.
\bibitem{ProjectCluster} P versus NP project, \emph{SLICECLUSTERMULT1}, 18 October 2026.
\end{thebibliography}
\end{document}
% FINAL_VERIFICATION_SENTINEL::MAIN::NAECOMPANION1::2026-10-18_1655::P_VS_NP_OPEN